\begin{helper}[Natural fiber mean scale]
For every \(\varepsilon>0\), if
\(m(n)=\noderef{TwoBiteNaturalReducedVertexCount}(n)\), then eventually
\[
  \noderef{WithinRelativeError}\!\left(
    \varepsilon,\;(\log n)^2,\;\frac{n}{m(n)}
  \right).
\]
Equivalently, the fixed-fiber hypergeometric mean \(n/m(n)\) is eventually
within relative error \(\varepsilon\) of \((\log n)^2\).
\end{helper}
\begin{proof}
Write \(L=\log n\), \(M=m(n)\), and
\(x=n/L^2\).  By \(\noderef{TwoBiteNaturalReducedVertexCount}\),
\(\noderef{TwoBiteReducedVertexCount}\), and \(\noderef{TwoBiteStretch}\),
we have \(M=\lfloor x\rfloor\).  Thus, for all sufficiently large \(n\) with
\(L>0\),
\[
  M\le x<M+1 .
\]
The existing scale lower bound
\(\noderef{TwoBiteNaturalReducedVertexRatioEventually}\) gives eventually
\[
  n\le 2L^2M .
\]
Since \(L^2=o(n)\), for the fixed \(\varepsilon>0\) we also eventually have
\[
  (2/\varepsilon)L^2\le n .
\]
Combining the last two inequalities and dividing by \(2L^2>0\) gives
\[
  1/\varepsilon\le M .
\]
In particular \(M>0\) and \(1/M\le\varepsilon\).

Now \(M\le x\) implies \(n/M\ge L^2\), while \(x<M+1\) gives
\[
  \frac{n}{M}-L^2
  =
  L^2\frac{x-M}{M}
  \le L^2\frac{1}{M}
  \le \varepsilon L^2 .
\]
The difference is nonnegative, so this is exactly the asserted
\(\noderef{WithinRelativeError}\) bound.
\end{proof}
